\documentclass[conference]{IEEEtran}
\usepackage[utf8]{inputenc}
\usepackage[T1]{fontenc}
\usepackage{cite}
\usepackage{amsmath,amssymb}
\usepackage{graphicx}
\usepackage{booktabs}
\usepackage{url}

\title{Verifier-Guided Evaluation of LLM-Generated Network Configurations: a Paired Study with Shared Initial Candidates}

\author{
\IEEEauthorblockN{Autonomous AI Researcher}
\IEEEauthorblockA{CNSM 2026 Agentic AI Researcher Track}
}

\begin{document}

\maketitle

\begin{abstract}
We preregistered and executed a paired confirmatory study (CAND-2026-001) that measured whether a verifier-triggered guarded pipeline (generate \ensuremath{\rightarrow} validate \ensuremath{\rightarrow} repair, <=1 repair) reduces deterministic-invariant violations relative to single-shot initial generation for intent\ensuremath{\rightarrow}configuration NetOps tasks. Using a deterministic synthetic task generator and a deterministic network verifier, we ran 200 preregistered tasks under shared\_initial\_candidate semantics with the hosted model gpt-5-mini (execution/model\_configuration.json records temperature = null/unset). Baseline configurations passed the verifier in 199/200 tasks (99.5\%); the guarded pipeline passed 200/200 (100\%). Paired absolute risk difference (guarded - baseline) = 0.005; contingency counts n11=199, n10=1, n01=0, n00=0; exact two-sided McNemar p = 1.0; paired-bootstrap 95\% percentile CI = [0.00, 0.015] (10,000 resamples, seed 1729). The single guarded improvement arose from one verifier-triggered repair that corrected a multi-step ordering violation. We report preregistration, full execution accounting (201 model calls: 200 shared initial + 1 repair), paired analysis, representative diagnostic artifacts, and bounded limitations grounded in the archived execution bundle.
\end{abstract}

\section{Introduction}
Generative large language models (LLMs) are increasingly proposed to translate high-level operator intent into device configurations for network operations (NetOps). Operational correctness for such workflows requires generated configurations to satisfy deterministic invariants (reachability, policy compliance, workflow ordering, and group-level safety). Prior prototype systems and benchmarks illustrate feasibility but leave open questions about how much deterministic verification and bounded repair alter acceptance rates when the identical initial sample is reused across comparison conditions.

This paper reports a preregistered paired experiment (CAND-2026-001) designed to isolate the marginal effect of a verifier-triggered automated repair on an identical generated candidate. Under shared\_initial\_candidate semantics the same single initial generation is deterministically reused by both baseline and guarded episodes; any paired difference therefore arises solely from the permitted validator-triggered repair. The primary estimand is the paired difference in per-task binary acceptance by a deterministic verifier (paired\_success\_rate\_difference\_guarded\_minus\_baseline). Because shared\_initial\_candidate semantics were enforced, this estimand measures a repair-only effect and does not capture potential benefits from independent guarded first-pass generations, constrained prompts, or multi-sample selection.

Contributions. (1) A preregistered, fully archived paired experiment executing 200 intent\ensuremath{\rightarrow}configuration tasks with a deterministic verifier and a hosted LLM (gpt-5-mini) under shared\_initial\_candidate semantics. (2) Complete execution accounting and deterministic reconciliation showing 200 shared initial generations and 1 repair call (201 model calls total). (3) Artifact-grounded confirmatory analysis reporting baseline pass 199/200 and guarded pass 200/200 with contingency counts (n11=199,n10=1,n01=0,n00=0), paired risk difference 0.005, exact two-sided McNemar p = 1.0, and bootstrap 95\% CI = [0.00,0.015] (10,000 resamples, seed 1729). (4) A localized failure diagnosis for the single repaired pair using archived validator traces. (5) Detailed reproducibility guidance and bounded limitations tied to archived artifacts and reconciliation records.

\section{Related Work}
We conducted a fresh literature investigation and verified records covering intent\ensuremath{\rightarrow}configuration systems, generate--validate--repair patterns, and agentic-AI safety discussions relevant to NetOps. Prior intent-to-configuration prototypes and task-bench evaluations demonstrate that LLMs can synthesize device configurations from operator intents and that verification is widely advocated as a guardrail [1],[2],[3]. Generate\ensuremath{\rightarrow}validate\ensuremath{\rightarrow}repair has empirical support in program repair and synthesis-with-tests domains and motivates guarded NetOps pipelines, but those demonstrations are primarily in software-engineering contexts rather than artifact-backed NetOps verifier studies [4],[5],[6].

Survey and safety literature on agentic AI documents substantive failure modes (hallucination/confabulation) and argues for verifier gates and provenance; this conceptual grounding motivated our guarded design and strict oracle definition but does not supply the verifier-backed NetOps empirical outcomes reported here [7],[8]. Synthesis. The supplied, verified records indicate active prototyping and conceptual advocacy for verifier-augmented pipelines but leave an empirical gap: few reproducible paired experiments quantify the marginal repair effect when the identical initial candidate is reused. Our paired shared\_initial\_candidate design directly addresses that gap.

\section{Methodology}
Preregistration and design freeze. The study design and analysis plan were frozen in preregistration CAND-2026-001 prior to any model calls. The preregistration fixed: (a) a paired binary estimand (paired\_success\_rate\_difference\_guarded\_minus\_baseline) comparing per-task binary verifier acceptance under baseline (no repair) versus guarded (verifier may trigger <=1 repair) conditions; (b) deterministic task selection for N=200 tasks; (c) generation\_semantics = shared\_initial\_candidate (one sampled initial generation per task deterministically reused by both episodes); (d) a maximum of one automated repair call per task; (e) primary failed-call treatment complete\_pair\_primary and a sensitivity treatment that counts persistent unparsables as failures; and (f) the confirmatory test (exact two-sided McNemar via exact binomial tail) plus a paired-bootstrap procedure for a percentile CI. The preregistration artifact (immutable) is referenced in the execution manifest.

Deterministic task sampling. Tasks were produced by a deterministic task generator (deterministic\_netops\_task\_generator\_v5) and selected by a fixed index rule documented in the preregistration. The selection rule is deterministic across the 200-task bank (the preregistration specifies the cycle-and-offset index mapping used to pick each task index); this ensures reproducible task composition and eliminates post-hoc selection. Each manifest encodes: initial\_state, expected\_final\_state, required\_sequence (ordered command list), allowed\_touched\_fields, workflow\_policies (group-level constraints, minimum-active rules), and difficulty metadata (changed\_interface\_count, dependency\_rule\_count, required\_command\_count). Difficulty labels used in the study include very\_hard and extreme; representative required\_command\_count values ranged from 3 up to 9 in the archived sample.

Shared\_initial\_candidate semantics and sampling notes. The adapter contract enforced shared\_initial\_candidate semantics: for each pair the exact same initial generation (a single sampled candidate) was provided to both the baseline and guarded episodes. Consequently, any observed paired difference can only arise from the action of the deterministic verifier and the single permitted repair; it cannot arise from independent first-pass guarded generations, different sampling calls, prompt-engineering differences, or selection among multiple generated candidates. The model-configuration record declares the requested model as gpt-5-mini and records temperature = null (unset). We explicitly note that temperature = null/unset in the configuration file is not evidence of deterministic sampling and does not imply temperature = 0; nondeterministic hosted-model sampling therefore remains a possible background property even though the experiment used a single initial sample per pair by design.

Generation, sanitizer, and repair policy. For each task the pipeline executed a deterministic formatting sanitizer (automatic\_syntax\_sanitizer\_v1) to correct minor syntactic normalization issues before verifier parsing; this sanitizer is deterministic and is applied prior to scoring per the preregistration. After the shared initial generation and sanitizer, the deterministic verifier parsed the candidate. If the verifier reported any invariant violations, the guarded logic was permitted to invoke at most one automated repair call to the hosted model with a repair prompt constructed from the validator trace (this repair policy and prompt format were specified in the preregistration). Repair calls were counted against the maximum\_model\_calls budget; the execution manifest reports exactly one repair call across all pairs in the run. The pipeline's missingness plan treated persistent unparsable outputs (after sanitizer and an allowed repair attempt) as scorable failures in the sensitivity analysis; in the primary analysis pairs with complete episodes remained eligible per the preregistered complete\_pair\_primary rule.

Deterministic verifier as the oracle. The deterministic verifier (deterministic\_netops\_validator\_v5) is the study's operational oracle. The verifier implements: syntactic parsing into discrete commands; enforcement of allowed\_fields and allowed\_touched\_fields per task manifest; comparison of parsed command sequences against the manifest's required\_sequence ordering; evaluation of group-level invariants (e.g., final\_group\_constraints and minimum\_active\_in\_groups); and a final boolean full\_intent\_constraint\_validation = true only if all encoded constraints are satisfied. The verifier emits structured tracer output for each episode that includes parsed\_command\_count, required\_command\_count, observed\_touched\_field\_count, normalized\_configuration, validation\_before (violation codes and counts), validation\_after (post-repair state and validity flag), validator\_id and version, and a human-readable description of detected invariant breaches. These tracer fields are the authoritative evidence for per-episode scoring in the analysis.

Paired execution bookkeeping and provider accounting. The preregistered planned episodes were 400 (200 pairs \ensuremath{\times} 2 conditions). The run completed all planned episodes without terminal provider failures. Provider accounting recorded hosted-model-call events and stage counts consistent with preregistration: the stage counts indicated 200 shared initial generation events and 1 repair event, for a total of 201 model calls. Cache and reuse checks verified that the shared initial candidate was reused per pair and that no cross-condition cache-reuse introduced additional independent generation calls. All runtime provider metadata (call counts, cache hits/misses, stage counts) were reconciled deterministically in the execution reconciliation artifact.

Statistical procedures --- confirmatory and CI computation. The primary confirmatory test is the exact two-sided McNemar test computed from discordant pair counts n10 and n01 (where n10 = count of pairs with guarded=1 \& baseline=0, and n01 = count with guarded=0 \& baseline=1). We implemented the exact two-sided McNemar p-value using the exact binomial-tail formulation: let n = n10 + n01 and k = n10; the two-sided p-value is p = 2 * min\{ BinomCDF(k; n, 0.5), 1 - BinomCDF(k - 1; n, 0.5) \} with the conventional handling when k = 0. This exact-binomial realization is documented in the analysis log and reproduces the p = 1.0 reported for the observed discordants (n = 1, k = 1). For the paired absolute risk difference we computed a paired-bootstrap 95\% percentile confidence interval with 10,000 resamples and bootstrap\_seed = 1729; bootstrap parameters and seed are recorded in the analysis log.

Missingness, exclusions, and contamination controls. The preregistration defined deterministic exclusion criteria prior to model calls (e.g., token-normalized identity to public vendor-source templates to avoid leakage) and required logging of any excluded tasks. The run recorded no preregistered exclusions. The sanitizer and missingness policy were applied deterministically; persistent unparsables did not occur in this execution. Contamination checks included token-level similarity scans between generated candidates and vendor-template sources and a contamination summary that would flag unusually high overlaps; no contamination flags were raised in this run.

Reproducibility and artifact grounding. All implementation details required to reproduce the methodology are archived in the execution bundle: the immutable preregistration, the deterministic task generator artifacts and manifests (task bank with required\_sequence and workflow\_policies), the model-configuration record (including declared model and the explicit temperature = null/unset entry), the verifier implementation and structured tracer outputs, provider-call accounting and deterministic reconciliation, the sanitizer and repair-prompt templates, and the analysis-executor configuration (exact McNemar code path and bootstrap seed/resample settings). The verifier's tracer fields provide the operational link between generated text and the binary scoring used in the confirmatory analysis. Re-running the archived paired analysis executor on the raw results reproduces the contingency counts and CI parameters reported in the paper.

\section{Results}
Execution accounting and provider audit. The execution manifest records start and end timestamps (UTC) and provider-call accounting: the run started at 2026-09-01T16:03:05.491039+00:00 and completed at 2026-09-01T16:32:13.298536+00:00 (\ensuremath{\approx}29 minutes wall-clock). Planned\_episode\_count = 400 and completed\_episode\_count = 400; failed\_episode\_count = 0. Provider-call audit (deterministic\_reconciliation.json) reports hosted\_model\_call\_event\_count = 201 with stage\_counts = \{shared\_initial\_generation: 200, repair: 1\}. These counts match the preregistered shared\_initial\_candidate semantics and the single observed repair.

Primary confirmatory outcomes (archived CSV artifacts). analysis/condition\_summary.csv shows: Baseline successes = 199, failures = 1 (success\_rate = 0.995); Guarded successes = 200, failures = 0 (success\_rate = 1.000). The paired contingency table (analysis/paired\_contingency\_table.csv) records n11 = 199, n10 = 1, n01 = 0, n00 = 0. The paired absolute risk difference (guarded - baseline) = 0.005. Exact two-sided McNemar p = 1.0 as computed above; paired-bootstrap percentile 95\% CI (10,000 resamples, seed 1729) = [0.00, 0.015]. All numeric values are present in analysis/*.csv and analysis/results.json and reproduced by the archived paired\_binary\_analysis\_v1 executor.

Representative repaired case and diagnostics. The guarded pipeline invoked exactly one automated repair. Validator traces for the repaired pair indicate an extreme composite task (required\_command\_count = 9, changed\_interface\_count = 3) where the shared initial candidate violated the task's required\_sequence ordering constraint. The repair applied a localized reordering that aligned the command sequence to required\_sequence and yielded validation\_after.valid = true. Per-episode validator\_trace fields (validation\_before.violation\_codes, parsed\_command\_count, required\_command\_count, normalized\_configuration) document both the pre-repair violation and the post-repair correction. These archived traces (execution/scoring/task-00000X-*.json) are the evidence for the localized repair effect and support the preregistered expectation that multi-device sequencing errors concentrate in higher-difficulty tasks.

Missingness, exclusions, and contamination. The run recorded no preregistered exclusions and no persistent unparsables; analysis/missingness\_summary.csv shows complete\_pairs = 200 and zeros for other missingness categories. analysis/contamination\_summary.csv is empty (no flagged contamination). Deterministic\_reconciliation.json marks all checks passed and pair\_accounting\_consistent = true. Provider-call audit shows cache\_miss\_count = 201 and cache\_hit\_count = 0.

Exploratory diagnostics by difficulty metadata. Representative task manifests illustrate required\_command\_count values from 3 to 9 and difficulty labels very\_hard or extreme. The single baseline failure is associated with an extreme pattern and high required\_command\_count, consistent with the preregistered exploratory question that multi-device sequencing and dependency complexity dominate residual failures. analysis/paired\_difference\_figure.svg (archived) visualizes condition rates and the bootstrap distribution used for the CI calculation.

\section{Discussion}
Interpretation under shared\_initial\_candidate semantics. Because baseline and guarded conditions reuse the identical sampled candidate per pair, observed differences isolate the marginal effect of permitting an automated verifier-triggered repair. The experiment therefore quantifies a repair-only effect and does not evaluate other guarded variants that would permit independent guarded first-pass generations or constrained prompting. We reiterate this estimand interpretation explicitly so readers do not generalize to independent-generation contrasts without new experiments.

Statistical regime, power, and limitations of inference. Our preregistered power rationale targeted detection of moderate absolute paired increases (\textasciitilde{}10--15 percentage points) at N = 200 under mid-range baseline assumptions (\textasciitilde{}45--65\%). The realized baseline success rate (199/200 = 99.5\%) places the analysis in a sparse-discordant regime where the exact McNemar p-value gives limited directional evidence (p = 1.0 for our observed discordants). The paired-bootstrap CI [0.00, 0.015] provides a bounded estimate of plausible effect sizes compatible with the data and shows the dataset is consistent with at most small absolute improvements under our repair-only design. This underscores that unexpectedly high baseline performance can substantially reduce realized inferential sensitivity for small marginal effects. Practically, when baseline acceptance is already near one, corrective repairs can still be operationally useful for rare but high-consequence failures; however, the present data quantify only that such rare corrections occurred at least once in the preregistered sample and do not permit precise frequentist inference about very small effect magnitudes beyond the bootstrap bounds reported.

Mechanistic interpretation of the single repair. The archived validator\_trace shows the repaired case involved a multi-step ordering violation across multiple interfaces and required\_command\_count = 9. The repair used by the guarded pipeline enacted a localized reordering consistent with the task manifest's required\_sequence and produced validation\_after.valid = true. This pattern aligns with our exploratory hypothesis (H3) that multi-device sequencing and inter-device dependency rules concentrate residual failures. From a mechanistic standpoint, verifier-encoded ordering constraints are a concretely checkable invariant; when the initial candidate contains an expressible sequence breach, a bounded repair that reorders commands according to manifest constraints can correct the fault. That corrective potential depends directly on (a) the verifier's ability to detect an expressible invariant breach and (b) the repair routine's capacity to deterministically construct a corrected sequence without introducing new invariant violations.

Operational implications and bounded claims. Artifact-grounded evidence supports a narrow operational conclusion: in this synthetic deterministic task bank using gpt-5-mini and deterministic\_netops\_validator\_v5, a verifier-triggered bounded repair corrected a localized multi-step ordering fault in one extreme task. This is direct evidence that verifier-guided bounded repair can remedy certain ordering/dependency violations in generated configurations when those violations are expressible in the task manifest's invariants. It is not evidence that bounded repair will correct omissions only observable at dataplane runtime, or that guarded repair will address adversarially poisoned contexts; those remain open evaluation axes requiring live emulation, multi-model studies, and adversarial contamination experiments. Moreover, because our repair allowance was limited to a single automatic repair call per task, the observed effect quantifies what a minimally permissive guarded pipeline can accomplish; pipelines that permit multiple repair iterations, richer external-tool synthesis, or human-in-the-loop adjudication may achieve larger gains but were not evaluated here.

Threats to validity and generalization (expanded). Internal validity: the paired shared\_initial\_candidate design isolates repair effects but does not address prompt sensitivity, model sampling variability, or the potential for different initial candidates under guarded prompting strategies. The model\_configuration.json records temperature = null (unset), which is not evidence of deterministic sampling; nondeterminism in hosted-model generation remains a background condition that the shared\_initial\_candidate semantics mitigated for the paired comparison but did not eliminate as a systemic source of variance. Construct validity: the verifier enforces encoded syntactic and ordering invariants but cannot capture behaviors observable only through dataplane testing (e.g., runtime routing convergence or vendor-specific command side-effects). External validity: the experiment used a single model (gpt-5-mini) and a synthetic, deterministic task bank; results do not imply cross-model or production-device generality. Conclusion validity: with a single discordant pair the realized inferential leverage for small repair-only effects is weak; the bootstrap CI quantifies that uncertainty and highlights the value of larger samples or targeted sampling focusing on high-difficulty strata for future studies.

Practical considerations for guarded pipeline design. The findings articulate two practical design lessons grounded in archived artifacts: (1) deterministic verifiers that explicitly encode ordering and group-level invariants can serve as effective selectors for when automated repair is necessary; the verifier's binary valid flag cleanly triggers bounded repair without human intervention in our run. (2) When invariants are expressible and repairs are constrained (for example, limited to reordering or restricted field modifications), a single focused repair may suffice to recover validity for complex tasks. Both lessons depend on careful manifest design that encodes the relevant invariants and on repair primitives that are conservative (minimize unintended changes) and auditable.

Reproducibility and next evaluation axes. The archived execution and analysis artifacts provide a complete trace for reanalysis and extension: task manifests with required\_sequence encodings, per-episode validator\_trace outputs, model-call accounting, and bootstrap parameters (seed = 1729, resamples = 10000). To better estimate small repair-only effects or to evaluate alternative guarded semantics, subsequent preregistered runs should consider (a) stratified oversampling of extreme/high-complexity tasks to increase discordant counts, (b) multi-model comparisons, and (c) relaxing shared\_initial\_candidate semantics to measure independent guarded first-pass benefits. Each such axis requires careful preregistration to avoid post-hoc selection and to preserve interpretability of paired or unpaired estimands.

\section{Limitations}
\begin{itemize}
\item Synthetic deterministic task bank and verifier: results reflect correctness under the chosen synthetic intent templates and deterministic verifier and do not guarantee similar performance on live heterogeneous production devices or telemetry.
\item Single-model experiment (gpt-5-mini): findings do not demonstrate multi-model generality.
\item Shared\_initial\_candidate semantics: baseline and guarded conditions began from the same initial generation; the study measures verifier/repair effects only and does not evaluate independent first-pass generation contrasts.
\item Limited adversarial/contamination testing: the run did not include systematic adversarial retrieval or poisoned-context attacks; such threats remain an open evaluation axis.
\item Oracle coverage: the deterministic verifier enforces encoded invariants but may omit runtime behaviors and device-specific semantics observable only through live execution; external validity to live deployments is therefore limited.
\end{itemize}

\section{Conclusion}
We executed a preregistered paired experiment measuring the marginal effect of a verifier-triggered repair under shared\_initial\_candidate semantics for LLM-generated network configurations. In 200 synthetic deterministic tasks with gpt-5-mini and deterministic\_netops\_validator\_v5, baseline acceptance was 199/200 and guarded acceptance was 200/200; a single automated repair corrected a multi-device ordering violation. Paired absolute risk difference = 0.005 with paired-bootstrap 95\% CI [0.00, 0.015] (10,000 resamples, seed 1729) and exact two-sided McNemar p = 1.0. Artifact-grounded evidence therefore supports a constrained conclusion: verifier-guarded bounded repair can correct localized ordering faults in this synthetic verifier-backed setting. Broader operational claims require additional experiments: multi-model studies, independent-first-pass guarded semantics, adversarial contamination testing, and live-device or dataplane emulation to capture runtime semantics not modeled by the deterministic verifier.

\section*{Disclosure Statement}
Disclosure: This run implemented preregistered study CAND-2026-001 and was executed autonomously after the autonomy lock. Declared model: gpt-5-mini (execution/model\_configuration.json records temperature = null/unset; null/unset is not evidence of deterministic sampling or temperature = 0). Orchestration adapter: hosted\_netops\_gvr\_v1. Exact initial master prompt archived at provenance/master\_prompt.txt (SHA-256 1872df1e\allowbreak{}1805d2d9\allowbreak{}6940456c\allowbreak{}a016bd66\allowbreak{}5d1d5196\allowbreak{}add77f5a\allowbreak{}cdf1582b\allowbreak{}b39b15ba). Preregistration fixed generation\_semantics = shared\_initial\_candidate (one sampled initial generation per task was deterministically reused by baseline and guarded episodes) and allowed at most one automated repair call per task. Model-call accounting in the archived execution manifest reflects 200 shared initial generation calls and 1 repair call (201 model calls total). The deterministic verifier used was deterministic\_netops\_validator\_v5. Humans selected the broad topic family and constrained the initial master prompt prior to the autonomy lock as permitted by the track; no human manually edited technical manuscript content after the autonomy lock. All provider traces, verifier logs, task manifests, model-configuration records, and analysis artifacts are archived in the execution bundle and indexed by the execution manifest.

\begin{thebibliography}{99}
\bibitem{ref1} Changjie Wang, Mariano Scazzariello, Alireza Farshin, Simone Ferlin, Dejan Kostić, Marco Chiesa, "NetConfEval: Can LLMs Facilitate Network Configuration?", 2024, 10.1145/3656296
\bibitem{ref2} Ryan Beckett, Aarti Gupta, Ratul Mahajan, David Walker, "A General Approach to Network Configuration Verification", 2017, 10.1145/3098822.3098834
\bibitem{ref3} Jack Gallifant, Majid Afshar, Saleem Ameen, Yindalon Aphinyanaphongs, Shan Chen, Giovanni Cacciamani, Dina Demner‐Fushman, Dmitriy Dligach, Roxana Daneshjou, Chrystinne Oliveira Fernandes, Lasse Hyldig Hansen, Adam Landman, Lisa Soleymani Lehmann, Liam G. McCoy, Timothy A. Miller, Amy C. Moreno, Nikolaj Munch, David Restrepo, Guergana Savova, Renato Umeton, Judy Wawira Gichoya, Professor Gary S. Collins, Karel G.M. Moons, Leo Anthony Celi, Danielle S. Bitterman, "The TRIPOD-LLM reporting guideline for studies using large language models", 2025, 10.1038/s41591-024-03425-5
\bibitem{ref4} Yunhe Li, "Test-in-the-loop LLM Repair: Verifiable Automated Program Repair on QuixBugs with a "Failing Test \ensuremath{\rightarrow} Patch \ensuremath{\rightarrow} Regression Test" Loop", 2024, 10.69987/jacs.2024.40206
\bibitem{ref5} Wayne Xin Zhao, Kun Zhou, Junyi Li, Tianyi Tang, Zican Dong, Yupeng Hou, Beichen Zhang, Yingqian Min, Junjie Zhang, Peiyu Liu, Xiaolei Wang, Yifan Du, Yushuo Chen, Yushuo Chen, Zhipeng Chen, Jinhao Jiang, Ruiyang Ren, Yifan Li, Xinyu Tang, Peiyu Liu, Yiwen Hu, Jian‐Yun Nie, Ji-Rong Wen, "A Survey of Large Language Models", 2026, 10.1007/s11704-026-60308-3
\bibitem{ref6} Haoran Dong, "Tool-Augmented Large Language Model Agents for Analysis: A Focused Study", 2026, 10.2139/ssrn.6647800
\end{thebibliography}

\end{document}
